\documentclass[11pt]{article}
\usepackage[a4paper,margin=1in]{geometry}
\usepackage{amsmath,amssymb}
\usepackage[hidelinks]{hyperref}
\usepackage{enumitem}
\usepackage{microtype}
\setlength{\parskip}{0.45em}
\setlength{\parindent}{0pt}

\title{SCWF Revision Summary}
\author{}
\date{March 11, 2026}

\begin{document}
\maketitle

\section*{What changed}

This revision keeps the self-consistent wavelet-frame estimator but removes three
major sources of ambiguity: spectral normalization, magnetic-unit ambiguity, and the
misinterpretation of the anisotropy buckets.

\begin{itemize}[leftmargin=1.6em,itemsep=0.35em]
\item The second-order spectral estimate now uses the positive-frequency discrete
response-energy integral of the exact FFT response used by the code.
\item The magnetic families are exported explicitly as \texttt{B\_nT} and
\texttt{B\_vel}. The compatibility family \texttt{B} still follows the user-selected
output units.
\item \texttt{Sfuncs} is retained only as a compatibility name. Its physically correct
interpretation is a scale-normalized wavelet-moment surrogate.
\item The local basis defines the angle bins, but the code averages the full vector
coefficient magnitude in those bins. It does not estimate basis-component spectra.
\end{itemize}

\section*{Defensible observables}

The revised package should be described in terms of the following observables.

For any coefficient family $W_j$ and conditioned subset $\mathcal{C}_j$,
\begin{equation}
 M_q^{(j,\mathcal{C})} = \langle |W_j|^q \rangle_{\mathcal{C}_j}
\end{equation}
is the raw wavelet moment, while
\begin{equation}
 \widetilde{M}_q^{(j,\mathcal{C})} = M_q^{(j,\mathcal{C})} / \tau_j^{q/2}
\end{equation}
is the saved scale-normalized surrogate.

For $q=2$, the conditional band-averaged PSD estimate is
\begin{equation}
 \widehat{P}^{(j,\mathcal{C})} = \frac{\langle |W_j|^2 \rangle_{\mathcal{C}_j}}{A_j^{(\mathrm{disc})}},
\end{equation}
where $A_j^{(\mathrm{disc})}$ is the positive-frequency discrete response-energy
integral of the exact implemented response.

\section*{What remains outside the scope of the method}

The package does not compute exact finite increments, strict two-point structure
functions, or basis-component Fourier spectra. The estimator should be presented as a
conditional wavelet-band analysis in a local anisotropy frame.

\section*{Practical reading guide}

\begin{itemize}[leftmargin=1.6em,itemsep=0.35em]
\item Use \texttt{WaveletMoments} when you want the raw coefficient moments.
\item Use \texttt{ScaleNormalizedWaveletMoments} when you want the higher-order
surrogate with the same field dimensions as an increment moment.
\item Use \texttt{Spectra} or \texttt{ConditionalBandAveragedPSD} when you want the
second-order conditional band-power estimate.
\item Check the metadata before reading the compatibility family \texttt{B}; its units
still depend on \texttt{return\_B\_in\_vel\_units}.
\end{itemize}

\end{document}
